\section{Table/Problem File Relationship}\label{sec:table-problem-relationship}

Information in the problem file tells TAOS which tables to use for each part of the
trajectory. Trajectories are divided into segments, and different tables can be used in
each segment. For example, the line

\begin{lstlisting}[style=taosmanual]
*aero ca=(stage1-ca-on) cn=(stage1-cn)
\end{lstlisting}

in the problem file tells TAOS to use a table called \tableid{stage1-ca-on} for the
axial-force coefficient and a table called \tableid{stage1-cn} for the normal-force
coefficient (\cref{sec:block-aero}). These table names must appear in the input table
files. When TAOS needs an axial-force coefficient, it retrieves the value from the
\tableid{stage1-ca-on} table.

Table names are given in \problemblock{*aero}, \problemblock{*prop},
\problemblock{*cg}, and \problemblock{*wind} data blocks in the problem file. They
can also be used in \problemblock{*define} data blocks when defining user-defined
output variables. All tables referenced in a problem file must appear in the table
files.

More than one table of each type can be given in a trajectory segment; the values
from each table are added together. For example, if two thrust tables are given, they
are evaluated separately and then added to obtain the total thrust force. Multiple
tables can also be used for a component buildup of the total vehicle lift and drag
forces.
